\section{Inverse-Calibration Problem}
\label{sec:inverse-problem}

For an observation vector $y\in\mathbb{R}^{20}$ of spot-normalized call and put
prices, the inverse problem is
\begin{equation}
    \hat{\bm{\theta}}(y)
    \in \arg\min_{\bm{\theta}\in\Theta}
    L\!\left(y, F(\bm{\theta};c)\right),
    \label{eq:inverse-calibration}
\end{equation}
where $F$ is the Double Heston pricer, $c$ contains maturity, rate, and carry
conditioning, and $\Theta$ enforces structural constraints and provisional bounds.
Traditional calibration minimizes a scale-normalized dollar-price residual with a
constrained optimizer. Neural models approximate $\hat{\bm{\theta}}(y)$ directly,
subject to their respective output maps and training losses.

The distinction between the forward map and its inverse is central. A small residual
\[
    \|y-F(\hat{\bm{\theta}};c)\|
\]
says that $\hat{\bm{\theta}}$ reproduces the observed surface under the selected norm.
It does not imply that $\hat{\bm{\theta}}$ equals the generating vector when multiple
near-equivalent vectors exist in $\Theta$. Consequently, all recovery statements in
this paper are tolerance-conditioned and use known synthetic truth.

Two error scales are used. Range-scaled errors divide each canonical coordinate by
the width of its reviewed hard-bound interval before computing an RMSE. Standardized
errors use training-split target statistics only. Both are reported without replacing
per-parameter diagnostics.
